\naslov{Metoda podpornih vektorjev}

Pri metodi podpornih vektorjev se ukvarjamo z dvojiško klasifikacijo.
Iščemo najširši rob, podan z afino preslikavo, da bo za pozitivne primere
veljalo $\beta^T x + \beta_0 \ge 1$, za negativne pa $\beta^T x + \beta_0 \le
-1$.
V obeh primerih postavimo ciljno spremenljivko na $y_i = \pm 1$, da je
$y_i(\beta^T x_i + \beta_0) \ge 1$.
Vsak učni primer nam da eno tako neenakost.
Če vzamemo dve točki na nasprotnem robu in upoštevamo, da je $\beta$ normala na
hiperravnini, v kateri se nahajata en pozitiven primer $x_+$ oz. en negativen
primer $x_-$, potem je širina pasu, ki ga hiperravnini definirata, enaka
\[
  \frac{\beta^T (x_+ - x_-)}{\norm{\beta}}
  = \frac{1 - \beta_0 - (-1 - \beta_0)}{\norm{\beta}}
  = \frac{2}{\norm{\beta}}.
\]
Želimo imeti čim širši pas, torej iščemo maksimum tega izraza; ekvivalentno
iščemo minimum $\pol \norm{\beta}^2$ pod pogojem, da veljajo zgoraj zapisane
neenakosti.

Minimum poiščemo z Lagrangeovimi multiplikatorji
\begin{gather*}
  L = \pol \norm{\beta}^2 - \sum_i \alpha_i (y_i(\beta^T \mb{x}_i + \beta_0) - 1), \\
  \partial_\beta L = \beta - \sum_i \alpha_i y_i \mb{x}_i, \\
  \partial_{\beta_0} L = - \sum_i \alpha_i y_i
\end{gather*}
Če enačimo zadnji enačbi z $0$ in vstavimo v $L$, dobimo
\[
  L = \sum_i \alpha_i - \pol \sum_{i,j} \alpha_i \alpha_j y_i y_j \mb{x}_i^T
  \mb{x}_j.
\]
Karush-Kuhn-Tuckerjevi pogoji nam dajo dodatne omejitve $\alpha_i \ge 0$ ter
\[
  \sum_i \alpha_i y_i = 0.
\]
Zahtevali smo $y_i (\beta^T \mb{x} + \beta_0) -1 \ge 0$, iz KKT robnih pogojev pa
dobimo $\alpha_i (y_i (\beta^T \mb{x} + \beta_0) -1) = 0$.
V primerih, ko je $\alpha_i \ne 0$, torej velja $y_i (\beta^T \mb{x} + \beta_0)
= 1$; to so ravno primeri na robu.
Pravimo jim \pojem{podporni vektorji}.
Napoved modela bo enaka
\[
  \hat{y}_0 = g\!\left( \sum_i \alpha_i y_i \mb{x}^T \mb{x}_0 + \beta_0 \right),
\]
kjer je $g$ indikatorska funkcija
\[
  g(x) =
  \begin{cases}
	1, & x \ge 0 \\
	0. & x < 0
  \end{cases}
\]

\vprasanje{Izpelji metodo podpornih vektorjev za linearno ločljive podatke.}

Če podatki niso linearno ločljivi, lahko v optimizacijsko funkcijo dodamo člen,
ki meri oddaljenost od pravilnega dela prostora, iščemo
\[
  \min_{\beta, \beta_0} \pol \norm{\beta}^2 + C \sum_i \xi_i,
\]
kjer je $\xi_i$ razdalja $i$-tega primera od roba, ki temu primeru pripada
(torej od pozitivnega roba, če je pozitiven, in od negativnega, če je
negativen), oziroma $0$, če je primer na pravilni strani roba.
Omejitve so sedaj oblike
\begin{gather*}
  y_i (\beta^T x_i + \beta_0) \ge 1 - \xi_i, \\
  \xi_i \ge 0.
\end{gather*}
Tudi za ta problem lahko najdemo optimalno rešitev s podobnim postopkom.
Parameter $C$ določa, kako pomembno je manjšanje števila podpornih vektorjev, ki
so sedaj lahko tudi znotraj roba (ali na drugi strani).
Večja vrednost $C$ pomeni, da bo model bolj ločeval razrede, s čimer se zniža
predsodek, a poveča varianca.

\vprasanje{Kaj narediš, če podatki v metodi podpornih vektorjev niso linearno
  ločljivi? Kaj je vloga parametra cene?}

Če imamo več kot dva razreda, lahko klasifikacijo naredimo na dva načina:
\begin{itemize}
\item Zgradimo model za vsak par možnih vrednosti v klasifikaciji
\item Zgradimo model za vsako možnost, ki ločuje podatke tega tipa od vseh ostalih
\end{itemize}

Če podatki niso linearni, jih transformiramo z neko funkcijo $\phi$, da
nelinearna odločitvena meja postane linearna.
Raje pa definiramo
\[
  K(u,v) = \sk{\phi(u), \phi(v)},
\]
ki nam dovoli, da se ne ubadamo s prehodom v transformiran prostor.
Jedrno funkcijo uporabimo namesto skalarnega produkta v predpisu za optimizacijo.
Na splošno lahko $\phi$ slika v poljuben Hilbertov prostor, zato so lahko jedra
konkretno komplicirana; če želimo, lahko $\phi$ recimo slika v
višjedimenzionalni prostor kot naš prostor podatkov.

\vprasanje{Kako modificiraš metodo podpornih vektorjev z jedrnimi funkcijami?}

% LocalWords:  hiperravnini višjedimenzionalni
